\appendix

\chapter{Полный перечень критериев оценки (C01--C25)}
\label{app:criteria}
\index{критерии оценки!полный перечень}

Ниже приведён перечень в редакции, согласованной с автоматическим скриптом оценки. Оцениваются прежде всего \textbf{изменения конкурсанта}; свойства учебных ЦР и Регулятора как отдельные критерии не входят. Каждый критерий --- до 3 баллов; номинальная сумма до \textbf{75}. Расшифровка уровней --- в приложении~\ref{app:criteria-rubric} и \path{docs/criteria\_rubric.md}.

\begin{longtable}{@{}p{0.12\textwidth}p{0.83\textwidth}@{}}
\toprule
\textbf{ID} & \textbf{Формулировка} \\
\midrule
\endfirsthead
\toprule
\textbf{ID} & \textbf{Формулировка} \\
\midrule
\endhead
C01 & Все тесты репозитория (включая тесты решения) завершаются успешно \\
C02 & Все тесты решения находятся в подкаталогах \texttt{src\_solution/tests/**} \\
C03 & Маркер \texttt{security} в \texttt{pytest.ini} и использование в тестах \\
C04 & Тесты журнала / \texttt{event\_log} \\
C05 & Пример \texttt{sga.json} \\
C06 & Примеры \texttt{SBOM\_TCB} / \texttt{SBOM\_OTHER} \\
C07 & Успешное выполнение сертификации решения (пакет + ответ Регулятора) \\
C08 & Сквозной автотест \textbf{ЦР--АБУ} (основной сценарий миссии), файл \texttt{tests/test\_e2e\_abu\_dm\_scenario.py} \\
C09 & Оформление кода в \texttt{src\_solution/} (flake8, PEP8) \\
C10 & Решение: журнал событий / \texttt{event\_log} в \texttt{src\_solution/} \\
C11 & Решение: зависимости в \texttt{src\_solution/} (\texttt{requirements.txt} / \texttt{pyproject.toml}); тяжёлые в ДВБ $\rightarrow$ 0, вынесение в \texttt{requirements-other.txt} при пустом \texttt{requirements.txt} $\rightarrow$ 3 \\
C12 & Тесты репозитория, \textbf{импортирующие} код из \texttt{src\_solution/} (AST) \\
C13 & Раздел тестов безопасности в отчёте решения в каталоге \texttt{src\_solution/docs/} \\
C14 & \texttt{numpy} в SBOM решения (\texttt{src\_solution/sbom/}): TCB vs OTHER (макс.\ 3) \\
C15 & Тесты репозитория: журнал \texttt{event\_log} и код в \texttt{src\_solution} \\
C16 & Покрытие ДВБ решения (\texttt{src\_solution/abu/tcb}) тестами \\
C17 & Отчёт конкурсанта в каталоге \texttt{src\_solution/docs/} \\
C18 & Решение: \textbf{security\_monitor}, \textbf{policies}; тесты политик (в т.\,ч. под \texttt{src\_solution/}) \\
C19 & Решение: изоляция доменов безопасности; монитор с контролем запросов и ответов \\
C20 & Стоимость сертификации --- место в рейтинге (жюри; в скрипте --- 0) \\
C21 & Экспертно: соответствие политик архитектуре АБУ (жюри; в скрипте --- 0) \\
C22 & Экспертно: полнота отчёта и воспроизводимость (жюри; в скрипте --- 0) \\
C23 & Размер домена ДВБ: оценка по максимальному LOC одного доверенного домена \\
C24 & Количество интерфейсов домена ДВБ: число разрешающих IPC-политик на домен \\
C25 & Наличие \texttt{security\_monitor} в решении, security-тесты и покрытие monitor-кода \\
\bottomrule
\end{longtable}

\chapter{Расшифровка критериев оценки (баллы 0--3)}
\label{app:criteria-rubric}
\index{критерии оценки!расшифровка}

Актуальная таблица уровней и путей к файлам поддерживается в репозитории: \path{docs/criteria\_rubric.md}. Номинальная сумма \textbf{raw} по автоматической части --- до \textbf{75}.

Каждый критерий \textbf{C01--C25} оценивается по уровням \textbf{0 / 1 / 2 / 3} (в т.\,ч. \textbf{C14} --- до 3 баллов за согласованное размещение \texttt{numpy} в SBOM\_TCB/SBOM\_OTHER). Экспертные \textbf{C20--C22} в скрипте дают \textbf{0}; жюри выставляет баллы вручную.

\section*{Автоматический скрипт}
Реализация: \path{scripts/evaluate\_contest\_score.py}. Особенности:
\begin{itemize}
  \item \textbf{C02:} все тесты решения должны находиться в \texttt{src\_solution/tests/**}.
  \item \textbf{C03:} маркер \texttt{security} в \texttt{pytest.ini} и использование в тестах.
  \item \textbf{C12:} подсчёт файлов в \texttt{tests/}, где AST содержит \texttt{import src\_solution}, а не подстрока в комментарии.
  \item \textbf{C13:} раздел тестов безопасности в отчёте решения в каталоге \texttt{src\_solution/docs/} с явными путями \texttt{src\_solution/...}.
  \item \textbf{C15:} тесты с импортами и \texttt{event\_log}, и \texttt{src\_solution}.
\end{itemize}

\section*{C01--C19 (кратко)}
Полный текст см.\ \path{docs/criteria\_rubric.md}. Уровни согласованы с версией скрипта в репозитории (релиз документации \textbf{1.7.7}).

\section*{C20: Стоимость сертификации (рейтинг)}
Выставляется \textbf{жюри} после анализа стоимости сертификации у всех участников (значения из логов/отчётов Регулятора).
\begin{itemize}
  \item \textbf{3} --- наименьшая стоимость среди конкурсантов.
  \item \textbf{2} --- второе место по стоимости.
  \item \textbf{1} --- третье место.
  \item \textbf{0} --- все остальные участники по этому критерию.
\end{itemize}
В автоматическом скрипте критерий даёт \textbf{0}; итог вносится в отчёт вручную.

\section*{C21: Экспертно --- соответствие политик архитектуре АБУ}
\begin{itemize}
  \item \textbf{0--3} --- жюри оценивает согласованность политик безопасности, реализованных в решении, с архитектурой АБУ и моделью угроз (в скрипте \textbf{0}).
\end{itemize}

\section*{C22: Экспертно --- полнота отчёта и воспроизводимость}
\begin{itemize}
  \item \textbf{0--3} --- жюри оценивает полноту отчёта, воспроизводимость команд, ясность аргументации (в скрипте \textbf{0}).
\end{itemize}

\section*{C23--C25: Дополнительные критерии по ДВБ и монитору}
\begin{itemize}
  \item \textbf{C23} --- размер доверенных доменов (максимальный LOC одного домена ДВБ).
  \item \textbf{C24} --- число разрешающих IPC-политик для доверенного домена.
  \item \textbf{C25} --- наличие \texttt{security\_monitor}, security-тестов и порогового покрытия monitor-кода.
\end{itemize}


\chapter{Пример отчёта участника (каталог \texttt{src\_solution/docs/} и таблица баллов)}
\label{app:example-report}
\index{отчёт!пример}

\textbf{Текстовый отчёт} по критерию \textbf{C17} ведётся в каталоге \path{src\_solution/docs/} (в репозитории --- заполняемый шаблон с обязательными разделами согласно приложению~\ref{app:criteria-rubric}). Ниже --- \textbf{пример калибровки} скрипта баллов для текущего дерева репозитория (в т.\,ч. \texttt{src\_solution/abu/tcb}, \texttt{src\_solution/sbom/}). Имена критериев и логика совпадают с \path{scripts/evaluate_contest_score.py}.

\begin{tabularx}{\textwidth}{@{}l X@{}}
\toprule
\textbf{Поле} & \textbf{Значение} \\
\midrule
Участник / команда & заготовка \texttt{src\_starting\_point} (прототип организаторов) \\
Коммит / тег & фактический коммит в репозитории \\
Дата прогона & по факту прогона \\
\bottomrule
\end{tabularx}

\vspace{0.8em}
\noindent\textbf{Автоматические проверки}

\begin{tabularx}{\textwidth}{@{}l X X@{}}
\toprule
Проверка & Результат & Примечание \\
\midrule
\texttt{make tests-all} & OK & типично \texttt{33 passed, 2 skipped} \\
Покрытие общее & из лога сертификации & поле \texttt{coverage\_percent} \\
Покрытие security & $\geq$ порога & \texttt{security\_coverage\_percent} \\
Сертификация & успех при подготовленном пакете & \texttt{certificate\_id} = SHA-256 архива \\
\bottomrule
\end{tabularx}

\vspace{0.8em}
\noindent\textbf{Фрагмент лога сертификации} (пример):

\begin{verbatim}
Результат сертификации: успешно
Стоимость (усл. ед.): <число из estimated_cost>
Сертификат (SHA-256 пакета): <64 hex>
\end{verbatim}

\vspace{0.8em}
\noindent\textbf{Баллы по критериям} (пример калибровки для заготовки без наполненного \texttt{src\_solution}):

\begin{longtable}{@{}c p{0.62\textwidth} r@{}}
\toprule
№ & Критерий & Баллы \\
\midrule
\endfirsthead
\toprule
№ & Критерий & Баллы \\
\midrule
\endhead
1 & C01: все тесты репозитория (включая тесты решения) завершаются успешно & 3.0 \\
2 & C02: тесты безопасности (\texttt{tests/security/} или \texttt{src\_starting\_point/tests/security/}) & 3.0 \\
3 & C03: маркер \texttt{security} в \texttt{pytest.ini} и использование в тестах & 3.0 \\
4 & C04: \texttt{event\_log} / журнал & 1.0 \\
5 & C05: пример \texttt{sga.json} & 3.0 \\
6 & C06: SBOM TCB / OTHER в примерах & 3.0 \\
7 & C07: \texttt{prepare\_certification\_bundle.sh} & 3.0 \\
8 & C08: сквозной автотест ЦР--АБУ & 3.0 \\
9 & C09: flake8 / PEP8 в \texttt{src\_solution} & 3.0 \\
10 & C10: \texttt{event\_log} в \texttt{src\_solution} & 1.0 \\
11 & C11: зависимости в \texttt{src\_solution} & 0.0 \\
12 & C12: импорты \texttt{src\_solution} в тестах & 1.0 \\
13 & C13: \texttt{security\_tests.md} и \texttt{src\_solution} & 2.0 \\
14 & C14: SBOM \texttt{numpy} (TCB vs OTHER) & 2.0 \\
15 & C15: тесты \texttt{event\_log} + \texttt{src\_solution} & 0.0 \\
16 & C16: покрытие \texttt{src\_solution/abu/tcb} & 3.0 \\
17 & C17: отчёт в каталоге \texttt{src\_solution/docs/} & 3.0 \\
18 & C18: \texttt{security\_monitor}, policies; тесты политик & 0.0 \\
19 & C19: изоляция доменов; монитор запрос/ответ & 0.0 \\
20 & C20: стоимость сертификации (рейтинг, жюри) & 0.0 \\
21 & C21: экспертно --- архитектура (жюри) & 0.0 \\
22 & C22: экспертно --- отчёт (жюри) & 0.0 \\
23 & C23: размер домена ДВБ (максимальный LOC) & 0.0 \\
24 & C24: количество интерфейсов домена ДВБ & 0.0 \\
25 & C25: monitor, security-тесты и покрытие monitor-кода & 0.0 \\
\midrule
\multicolumn{2}{@{}r}{\textbf{Сумма (raw):}} & \textbf{37.0 / 75} \\
\bottomrule
\end{longtable}

\vspace{0.5em}
\noindent\small Критерии C20--C22 (жюри) в автоматике 0; при отсутствии кода в \texttt{src\_solution/} соответствующие пункты дают 0.

\chapter{Регламент соревнований (полный текст)}
\label{app:regulations}
\index{регламент}

\section*{Термины}

\begin{tabularx}{\textwidth}{@{}l X@{}}
\toprule
\textbf{Термин} & \textbf{Определение} \\
\midrule
SGA & \textit{security goals and assumptions}; в коде и API; в текстах допустимо ЦПБ. \\
Регулятор & сертификация пакета АБУ. \\
ЦР & парк установок и миссии. \\
\bottomrule
\end{tabularx}

\section*{Последовательность действий конкурсантов}

\begin{enumerate}
  \item Развернуть окружение: \texttt{make install}, при необходимости Docker (\texttt{make docker-up}).
  \item Изучить \texttt{docs/context.md}, \texttt{docs/architecture.md}, \texttt{docs/certification\_process.md}.
  \item Разрабатывать решение в \texttt{src\_solution/} (или согласованный каталог), опираясь на \texttt{src\_starting\_point/}.
  \item Прогонять \texttt{make tests-all} и сертификацию (\texttt{make prepare-cert-bundle}, \texttt{make certify-abu} или API Регулятора).
  \item Заполнить отчёт о решении в каталоге \texttt{src\_solution/docs/} (критерий C17) и таблицу баллов по \texttt{docs/templates/evaluation\_report.md}.
  \item Подать работу: ссылка на \textbf{приватный} git-репозиторий; добавить пользователя организаторов с правами \textbf{разработчика} (GitHub: Collaborator с доступом на запись; GitLab: роль Developer или выше --- по инструкции жюри).
\end{enumerate}

\section*{Запрещённые приёмы}

\begin{itemize}
  \item Подмена или отключение автоматических тестов заготовки без обоснования в отчёте.
  \item Обход Регулятора (подделка сертификата, недоступный для проверки хэш).
  \item Публикация секретов, ключей, персональных данных в репозитории.
  \item Плагиат без ссылки на источник.
\end{itemize}

\section*{Критерии оценки (обобщение)}

25 критериев (C01--C25), каждый до 3 баллов; номинальная сумма до 75. Автоматическая часть: тесты (в т.\,ч. сквозной сценарий ЦР--АБУ), артефакты сдачи, эвристики по \texttt{src\_solution/} (включая C18--C19 и C23--C25). Экспертная часть: C20 (рейтинг по стоимости сертификации), C21--C22 --- жюри (в скрипте 0).

Подсказка по C19--C20: учебный пример изоляции ---\\
\path{references/traffic_light_demo/cyberimmunity-traffic-lights-example.ipynb}

\section*{Разрешение равенства баллов}

\begin{enumerate}
  \item Больший балл за автоматическую часть (в т.\,ч. C01--C19).
  \item Учёт C20--C22 (в т.\,ч. рейтинг по стоимости сертификации и экспертные критерии).
  \item Более ранняя успешная сдача.
  \item Меньше предупреждений/ошибок в отчёте сертификации.
  \item Жюри по чек-листу (порядок объявляется до старта).
\end{enumerate}

\section*{Прозрачность}

Единый шаблон отчёта; участник вправе апеллировать в оговорённые сроки, предоставляя логи и воспроизводимые команды.

\chapter{Шаблон отчёта и порядок апелляции}
\label{app:template-appeal}
\index{апелляция}
\index{шаблон отчёта}

\section*{Структура отчёта (шаблоны)}

\textbf{Текстовый отчёт} (критерий C17) --- файл в каталоге \path{src\_solution/docs/} в репозитории: архитектура, политики, сквозные и security-тесты, сертификация, диаграммы (см. приложение~\ref{app:criteria-rubric}, C17).

\textbf{Таблица баллов} для прозрачности --- \path{docs/templates/evaluation\_report.md}. Участник заполняет разделы:

\begin{tabularx}{\textwidth}{@{}c X@{}}
\toprule
\textbf{№} & \textbf{Раздел} \\
\midrule
1 & Автоматические проверки (таблица: проверка, результат, примечание) \\
2 & Лог сертификации (фрагмент stdout / API) \\
3 & Баллы по 25 критериям (0--3 за критерий), сумма / 75 \\
4 & Разрешение равенства баллов (при равенстве суммы) \\
5 & Апелляция: срок и контакт --- по объявлению жюри \\
\bottomrule
\end{tabularx}

\vspace{0.8em}
\noindent\textbf{Таблица критериев в отчёте} (формат строки):

\begin{tabularx}{\textwidth}{@{}c X c@{}}
\toprule
№ & Критерий & Баллы (0--3) \\
\midrule
1--25 & C01--C25 (см. приложение~\ref{app:criteria}) & заполнить \\
\bottomrule
\end{tabularx}

\section*{Порядок апелляции}

\begin{enumerate}
  \item Апелляция подаётся в рамках сроков, объявленных жюри для конкретного этапа олимпиады.
  \item Участник предоставляет артефакты прозрачности: логи прогонов, воспроизводимые команды (\texttt{make tests-all}, \texttt{make certify-abu}, и т.\,д.), ссылку на коммит.
  \item Контакт и форма подачи указываются в объявлении жюри; до публикации сроков --- ориентир: раздел «Прозрачность» регламента (приложение~\ref{app:regulations}).
\end{enumerate}
